\section{Empirical Strategy}
\label{sec:empirical}

\subsection{The Empirical Object: Discrimination of Cartel-Adjacent
Populations Under Coarsened Observability}
\label{sec:emp_estimand}

The primary empirical object is the discrimination accuracy of the
participation-based statistic against cartel-adjacency labels
under the data envelope the framework constrains us to: an award-
record layer that preserves participant identity, winner identity,
item code, and negotiated price, but not per-bidder bid amounts.
The discrimination object is what the framework's sufficient
ranking statistic predicts will dominate when only the award layer
is observable, and it is the object the validation evidence in
\S\ref{sec:cade} carries.

The conditional log-price association reported in
\S\ref{sec:results} is a secondary object: a descriptive
corroboration that the screening signal also leaves a price
imprint on the items it concentrates on, not a treatment-effect
target. Under the separating equilibrium of Online Appendix~A,
frequent-loser participation is generated by the cartel's
deployment problem; endogenous sorting into items where the
cartel deploys is constitutive of the screening object. The
screening question asks how informative the resulting
participation footprint is about the collusive environment when
only the award layer is observable; the relevant estimator
integrates over the distribution of items into which frequent
losers actually enter.

We therefore report two empirical objects, with the discrimination
object as the primary defense and the conditional association as
descriptive corroboration:

\paragraph{Primary object: discrimination AUC} the receiver-operating
characteristic AUC of the binary frequent-loser indicator and the
underlying continuous statistic $\log(1+\text{tenders\_count})$
against three cartel-adjacency labels: (i) cobidders inside the
always-loser stratum on a strict pre-$2020$ contemporaneous
benchmark with participation-stratified permutation null
(\S\ref{sec:cade_contemporaneous}); (ii) cobidders under temporal
holdout, train $2009$--$2016$ and test $2017$--$2019$
(\S\ref{sec:cade_prospective}); and (iii) the same target read
against the \citet{imhof2019detecting}--style bid-distribution
pipeline trained on the bid layer (\S\ref{sec:forensic}). The
discrimination object is what the framework's sufficient ranking
statistic predicts will dominate; the AUC numbers are the
empirical content the screening interpretation rests on.

\paragraph{Secondary object: broad-sample conditional association \texorpdfstring{($\beta$)}{(beta)}} the
within-cell conditional mean difference in log negotiated price
between tender-items that include at least one frequent-loser
participant and otherwise comparable items in the same product
code, year, and procuring unit. We estimate $\beta$ on
$\valSampleN$ tender-items and bracket it with four estimators
(OLS, cross-fit, CEM, IPW). It is descriptive corroboration that
the screening signal concentrates on items that also carry a
price imprint, not a causal target. We also report the
overlap-restricted association $\beta^{\text{ov}}$
(Table~\ref{tab:item_level_scope_match}), and the decomposition in
\S\ref{sec:results_overlap} shows that the broad-sample positive
imprint survives an unweighted overlap restriction at
$\valBetaUnwOverlap$ and reverses to $\valMatchOverlapCoef$ only
under ATT-reweighting that up-weights cells with relatively few
untreated counterfactuals.

\paragraph{Two design-based strategies that return null}
A sharp regression discontinuity at the convite/preg\~ao
statutory value caps is gated by the absence of bunching:
McCrary-style density tests around the $\valCapOldR$ cap return
a left/right ratio of $\valMcCraryRatio$ and a log-density
discontinuity of $\valMcCraryDisc$
(Table~\ref{tab:mccrary}), with the convite share dropping
cleanly from $\valPctSixnine$ below to $\valPctFifnine$ above.
Local first-stage power is too weak to support a fuzzy RDD at
standard bandwidths. A difference-in-differences design using
Decreto~$\valDecreto$'s $2018$ raise of the convite cap fails
parallel trends at the available pre-treatment window. We
report both failures openly. They do not constrain the
screening interpretation, which does not require a causal
identification of $\beta$, but they do constrain alternative
framings of the empirical object that would.

\subsection{Reduced-Form Specification}
\label{sec:emp_reduced}

The reduced-form specification is
\begin{equation}
  y_{igt} = \beta \cdot \text{losers}_{igt}
  + \mathbf{x}_{igt}' \boldsymbol{\delta}
  + \alpha_g + \lambda_t + \gamma_k + \varepsilon_{igt}
  \label{eq:ols_main}
\end{equation}
where $y_{igt}$ is the outcome for tender-item $i$ in item group
$g$ at time $t$ and procuring unit $k$; $\text{losers}_{igt}$
indicates the presence of at least one frequent-loser participant;
$\mathbf{x}_{igt}$ includes a convite indicator and any
specification-specific controls; $\alpha_g$, $\lambda_t$, and
$\gamma_k$ are item, year, and procuring-unit fixed effects; and
$\varepsilon_{igt}$ is the disturbance, clustered at the item
level. The headline outcome is $\log p_{\text{negotiated}}$;
auxiliary outcomes (number of bidders, number of bids, number of
non-frequent-loser bidders) are reported in
§\ref{sec:results_ols} and the appendix.

We estimate Equation~\eqref{eq:ols_main} under four
specifications. The general within-item-and-year baseline drops
the procuring-unit fixed effect; the within-procuring-unit
specification adds it. Two modal subsamples, preg\~ao-only and
convite-only, decompose the headline along the observability-regime
contrast set up in §\ref{sec:institutional}. The estimation
sample comprises $\valSampleN$ tender-items
(§\ref{sec:data_sample}).

Three complementary estimators bracket
Equation~\eqref{eq:ols_main}. Cross-fitting defines the
frequent-loser list on odd years, estimates the equation on
even years, and reverses the assignment, removing within-sample
classification noise. Coarsened exact matching (CEM) and
inverse-probability weighting (IPW) adjust for selection on
observables (year, procedure type, item group,
procuring-unit-size quartile). Under the screening framing all
three estimators bracket the broad-sample object $\beta$ rather
than a treatment-effect target; the four-estimator bracket
gives the conditional range $\valHeadlineRange$ used throughout
the paper (§\ref{sec:results_ols}). A leave-one-out
instrumental-variable specification uses the supply of frequent
losers in other procuring units in the same product market and
year as instrument; the IV coefficient is informative as a
direction-of-attenuation diagnostic, returning a magnitude
larger than OLS, the direction predicted when the binary
indicator coarsens the underlying continuous loss-intensity
primitive (§\ref{sec:mech_horse_race}). We do not interpret the
IV estimate as a causal magnitude.

Standard errors are clustered at the item level as the baseline
choice. Robustness to procuring-unit clustering and to two-way
\citet{cameron2011robust} clustering is reported in
§\ref{sec:rob_threshold}.

\subsection{Identification Audits as Screening-Value Diagnostics}
\label{sec:emp_audits}

The screening framing reorganizes the role of the four
identification audits reported in §\ref{sec:rob_identification}.
Each audit answers a distinct question about the screening
object rather than a question about the credibility of a causal
estimate.

The robustness values of \citet{cinelli2020making}
($RV_{q=1}=\valRVqOne$) and \citet{oster2019unobservable}
($\hat\delta=261.64$) bound the broad-sample association against
selection on observables already absorbed by item, year, and
procuring-unit fixed effects. They calibrate how much
unobserved-confounder variance would have to load on
observables to explain $\beta$ away. Under the screening framing
this is the bound on selection-on-observables explanations of
the screen's signal.

Strict-overlap matching delivers the second empirical object
$\beta^{\text{ov}}$
(Table~\ref{tab:item_level_scope_match}). Its role is not to
sharpen $\beta$ but to instantiate the alternative empirical
object whose contrast with $\beta$ identifies the screening-value
interpretation.

The leakage audit (Table~\ref{tab:leakage_audit}) decomposes the
in-sample item-level discrimination AUC into a structural
component and a tautology component, by holding cobidder firms
out of score construction within cross-validation folds and by
evaluating under temporal holdout. It is a generalization
diagnostic for the screening statistic, not an identification
test for $\beta$.

The sub-threshold-always-loser placebo
(Table~\ref{tab:iv_placebo}) confirms that the participation
primitive operates through the threshold-crossing population the
construct identifies, rather than through generic always-loser
supply. It is a specification diagnostic for the participation
primitive, not an identification test for $\beta$.

The four audits are reported in §\ref{sec:rob_identification}
and the identification scope they delineate is consolidated in
§\ref{sec:lim_identification}.

\subsection{Heterogeneity Across Detection Regimes}
\label{sec:emp_heterogeneity}

The framework of Online Appendix~A maps procuring-unit size to
the principal detection cost $\theta_k$ through the comparative
static $\partial m^{*}/\partial\theta_k<0$. We split the sample
into quartiles of procuring-unit size (annual contract volume
in BEC) and re-estimate Equation~\eqref{eq:ols_main} within
each quartile, with robustness to alternative size measures
(cumulative item count, headcount of distinct purchasers).
Under the screening framing the gradient reads the screening
signal's informativeness across detection regimes
(§\ref{sec:results_oversight}); we do not interpret it as
identification of an oversight channel. A second descriptive
heterogeneity specification, in winner-concentration and
repeated-pair density at the frequent-loser-firm level, is
reported in §\ref{sec:lim_construct} as scope of the screening
object rather than as a positive test of cartel-rotation
predictions.

\paragraph*{Summary}
The empirical strategy is organized around the discrimination
object first and the conditional association second. The
validation in \S\ref{sec:cade} carries the discrimination
evidence against three cartel-adjacency targets and the within-
stratum theory-validation bridge to the cover-bidder type. The
within-item price association in \S\ref{sec:results} reports
the broad-sample $\beta$, its overlap-restricted counterpart
$\beta^{\text{ov}}$, and the empirical decomposition of the
sign reversal as a reweighting phenomenon rather than an
overlap-dropping phenomenon. The architectural test in
\S\ref{sec:forensic} reads the screening statistic against bid-
distribution detectors trained on bid microdata. None of the
defended claims relies on a causal identification of $\beta$,
and the failed RDD/DiD designs reported above do not constrain
any of them.
